\begin{tomtat}
\begin{itemize}
  \item Khảo sát của Knab và cộng sự là một khảo sát tài liệu và nó không chạy
  thí nghiệm nào:
  không có mục kết quả, không dựng lại kỹ thuật nào, không có bộ dữ liệu riêng.
  Ba bảng của nó phân loại công trình đã công bố. Vì vậy mọi phát biểu ở chương
  này về việc LIME không ổn định là phát biểu của những bài báo mà khảo sát dẫn,
  thuật lại qua lời của khảo sát.
  \item Trục thứ nhất của phân loại là năm nhóm vấn đề: tính cục bộ, độ khớp của
  mô hình thay thế, khả năng diễn giải, tính ổn định
  (định nghĩa~\ref{def:ch10-tinh-on-dinh}) và hiệu năng. Năm nhóm không độc lập:
  khảo sát nêu rằng tăng tính cục bộ có thể làm hiệu năng xấu đi, còn hạ hiệu
  năng có thể làm tính ổn định tốt lên. Nhóm thứ hai được khảo sát gọi là
  \enquote{fidelity} và hỏi đầu ra của mô hình thay thế có bám đầu ra của mô hình
  gốc trên vùng đã lấy mẫu hay không. Sách gọi đó là độ khớp, giữ nó tách khỏi
  định nghĩa~\ref{def:ch01-do-trung-thuc} và tách khỏi đại lượng phản thực của
  chương~\ref{ch:khung-hop-nhat}.
  \item Trục thứ hai là bốn bước của quy trình: \tn{sinh đặc trưng}{feature
  generation}, \tn{sinh mẫu}{sample generation}, attribution đặc trưng, và
  \tn{biểu diễn lời giải thích}{explanation representation}. Bốn bước ấy là
  hình~\ref{fig:ch03-quy-trinh-lime} gộp thô hơn, và chỗ đứng của một biến thể
  trong bốn bước cho biết nó đã đổi mệnh đề nào về $f$.
  \item Bảng phân loại có 48 kỹ thuật: 23 nhắm vào độ khớp, 20 vào tính ổn định,
  13 vào khả năng diễn giải, 11 vào tính cục bộ, 3 vào hiệu năng, tổng cộng 70
  lượt gán vì 21 kỹ thuật mang từ hai nhóm trở lên. Đếm theo kỹ thuật thì 32
  trong 48, tức hai phần ba, nhắm vào độ khớp hoặc tính ổn định. Hiệu năng là
  nhóm duy nhất đo được mà không cần dụng cụ đo nào đang tranh cãi, và cũng là
  nhóm ít người làm nhất.
  \item Phần thảo luận của khảo sát ghi ba điều: một chuẩn thực hành cho nghiên
  cứu và đánh giá chưa được thiết lập, 50\% số kỹ thuật không có mã nguồn theo
  cách đọc mà mục~\ref{sec:ch10-khong-co-cach-so-sanh} nêu, và
  nhiều bài chỉ so bản sửa của mình với bản LIME chuẩn nên sau 48 công
  trình vẫn không có mệnh đề nào so hai biến thể với nhau.
  Hình~\ref{fig:ch10-vong-bien-the} vẽ vòng ấy: không bước nào trong vòng loại
  bỏ được một biến thể.
  \item Chữ \enquote{faithfulness} không xuất hiện lần nào trong phần thân bài, và
  \enquote{ground truth} cũng vậy. Bảng tính chất đánh giá 17 dòng của khảo sát mở
  đầu bằng \enquote{Correctness}, tức định nghĩa~\ref{def:ch01-do-trung-thuc} dưới
  cái tên thứ ba sau faithfulness và fidelity, và không dòng nào kèm công thức
  hay thủ tục tính. Một trong bốn hướng nghiên cứu mà khảo sát đề xuất là xây
  một khung đánh giá riêng cho LIME từ những chỉ số sẵn có; đọc cùng
  chương~\ref{ch:chi-so-khong-do-do-trung-thuc}, đó là xếp hạng 48 biến thể bằng
  một dụng cụ chưa ai kiểm định. Mối liên hệ ấy là của sách, không phải của
  khảo sát.
\end{itemize}
\end{tomtat}

\cauhoi
\begin{cauhoids}
  \item Mục~\ref{sec:ch03-cac-tham-so-tu-do} liệt kê bảy lựa chọn tự do của
  LIME, còn mục~\ref{sec:ch10-bon-buoc} xếp chúng vào ba trong bốn bước. Hãy
  chỉ ra bước nào không nhận lựa chọn nào trong bảy lựa chọn ấy, rồi cho biết
  vì sao vẫn có những biến thể can thiệp vào bước đó.
  \item Lấy $f(x_1,x_2) = x_1$ và một tập dữ liệu mà mọi điểm đều có
  $x_1 = x_2$. Hãy kiểm rằng $g_1 = x_1$ và $g_2 = x_2$ cùng khớp $f$ hoàn hảo
  trên tập ấy. Từ đó cho biết mỗi phát biểu sau đúng hay sai, và vì sao: hai lời
  giải thích khác nhau thì ít nhất một cái khớp sai; hai lời giải thích khác
  nhau thì chúng không cùng trả lời một câu hỏi về phụ thuộc; một phương pháp
  không ổn định thì con số độ khớp của nó không nói được gì về mô hình. Điều đó
  nói gì về 10 kỹ thuật mà bảng phân loại xếp vào cả nhóm tính ổn định lẫn nhóm
  độ khớp.
  \item Trong năm nhóm vấn đề, hiệu năng là nhóm có ít kỹ thuật nhất và cũng là
  nhóm dễ đo nhất. Hãy nêu hai cách giải thích khác nhau cho cặp sự kiện đó, một
  cách quy về giá trị khoa học của vấn đề và một cách quy về cách công trình
  được đánh giá, rồi cho biết dữ liệu nào phân biệt được hai cách giải thích ấy.
  \item Hình~\ref{fig:ch10-vong-bien-the} có hai mũi tên chấm đi vào ô câu hỏi.
  Giả sử một mũi tên được nối lại: toàn bộ 48 kỹ thuật đều công bố mã nguồn. Hãy
  cho biết vòng ấy có khép được không, và nếu chưa thì còn thiếu gì.
  \item Chữ \enquote{faithfulness} vắng mặt trong phần thân của khảo sát, và
  chương~\ref{ch:do-do-trung-thuc} đã ghi một trường hợp cùng loại. Hãy nêu một
  hệ quả cụ thể của việc hai nhánh nghiên cứu dùng hai bộ từ vựng không giao
  nhau cho cùng một tính chất, ở mức thao tác của người đi tìm tài liệu.
  \item Bảng tính chất của khảo sát mô tả \enquote{Correctness} bằng lời và không
  kèm thủ tục tính. Hãy chọn một họ chỉ số của
  mục~\ref{sec:ch08-ho-xoa-dan} và cho biết cần thêm những gì để biến câu mô tả
  ấy thành một thủ tục chạy được, rồi chỉ ra lựa chọn nào trong số đó là lựa
  chọn tự do theo nghĩa hình~\ref{fig:ch08-vong-lap-do}.
  \item Đọc \enquote{Which LIME should I trust? Concepts, Challenges, and
  Solutions} của Knab và cộng sự~\autocite{p22whichlime}, phần thảo luận. Hãy
  chọn một trong ba quan sát của phần ấy và thiết kế một phép đo cho nó trên
  một nhánh phương pháp khác trong sách, chẳng hạn nhánh attribution theo
  gradient của chương~\ref{ch:attribution-theo-gradient}, rồi cho biết phép đo
  ấy cần dữ liệu gì mà một khảo sát tài liệu không tự có.
\end{cauhoids}
